% YY12100B
% Einschätzungsblock
% 14.0
% 2013-11-30
:ref:`m-ersteinschaetzung`:ref:`m-einschaetzungsblock`%
\par\noindent
% \TableWide
% 	{<Titel>}%1
% 	{<Beschreibung>}%2
% 	{<Label>}%3
% 	{<Quelle>}%4
% 	{<Lizenz>}%5
% 	{<Tabelle>}%6
% 	{<Float>}%8
\TableWide{Einschätzungsblock}
{}
{tab:einschaetzungsblock}
{}
{}
{
\CorrectionItemizeBox\renewcommand{\itemhook}{%
%   \setlength{\topsep}{0pt}%
  \setlength{\itemsep}{0.004\baselineskip}
%   \setlength{\parskip}{\baselineskip}
  \setlength{\leftmargin}{1.4em}
  \setlength{\labelwidth}{0.9em}
}
\noindent\begin{minipage}[t]{\linewidth -
{2\fboxrule} - {2\fboxsep}}
\noindent\begin{minipage}[t]{\linewidth}
{\vspace{-2\baselineskip}%
\scriptsize
% \begin{tabulary}{\linewidth}[H]{cp{8em}p{8.5cm}}
\begin{tabularx}{\linewidth}[H]{p{1.5cm}XX}
% \toprule\addlinespace
&
\TabHeading{Beurteilung / Untersuchung}
&
\TabHeading{Typische Sofortmaßnahmen}
\\\addlinespace\toprule\addlinespace
\noindent\TabHeading{Szeneüberblick m. Schutz}
&
\begin{minipage}[t]{\linewidth}
\begin{itemize}
  -   Sicherheit, Selbstschutz; Patientenanzahl; 
  -   Umgebung, Ort,
Zeit, Gefahrenzonen, Umstände,
-   Weitere Resourcen erforderlich? Unfallmechanismus,
Großschaden
-   Lagemedlung erforderlich
-   \EE{Trauma?} \EE{Mechanismus?}
\end{itemize}
\end{minipage}
&
\noindent\color{Indigo}\begin{minipage}[t]{\linewidth}
\begin{itemize}
  -   GAS-Maßnahmen 
  -   Lagemeldung
  -   Schutzausrüstung
  -   Nachforderung weiterer Kräfte
\end{itemize}
\end{minipage}
\\\addlinespace\toprule\addlinespace
\noindent\TabHeading{Eindruck}
&
\begin{minipage}[t]{\linewidth}
\begin{itemize}
  -   Alter, Geschlecht, **Hautfarbe**, Gesichtsausdruck,
  **Haltung**, spontane Bewegungen,  Sprache
  -   Offensichtliche Verletzungen, Blutungen
  -   Sonstige **Auffälligkeiten**
\end{itemize}
\end{minipage}
&
\noindent\color{Indigo}\begin{minipage}[t]{\linewidth}
\begin{itemize}
  -   Manuelle HWS-Fixierung
  -   Stillung von starken Blutungen
  -   Bewegungsverbot
\end{itemize}
\end{minipage}
\\\addlinespace\midrule\addlinespace
\noindent\TabHeading{Bewusstsein}
&
\begin{minipage}[t]{\linewidth}
\begin{itemize}
  -   **Bewusstseinsgrad** (WASB)
\end{itemize}
\end{minipage}
&
\noindent\color{Indigo}\begin{minipage}[t]{\linewidth}
\begin{itemize}
  -   NA-Nachforderung
\end{itemize}
\end{minipage}
\\\addlinespace\midrule\addlinespace
\noindent\TabHeading{Hauptbeschwerde}
&
\begin{minipage}[t]{\linewidth}
\begin{itemize}
  -   **Berufungsgrund** und
  **Leitsymptome**
\end{itemize}
\end{minipage}
&
\noindent\color{Indigo}\begin{minipage}[t]{\linewidth}
\end{minipage}
\\\addlinespace\toprule\addlinespace
\noindent\T{A}

\TabHeading{Atemweg}
&
\begin{minipage}[t]{\linewidth}
\begin{itemize}
  -   *Inspektion der Atemwege (**Mund**, **Nase**)
  -   **Atemgeräusche**
\end{itemize}
\end{minipage}
&
\noindent\color{Indigo}\begin{minipage}[t]{\linewidth}
\begin{itemize}
  -   Absaugung, Fremdkörper entfernen
  -   Kopf überstrecken, Esmarch-Handgriff
  -   HWS-Immobilisation manuell oder mit Schiene
  -   \Z{Erweiterte Maßnahmen}
\end{itemize}
\end{minipage}
\\\addlinespace\midrule\addlinespace
\noindent\T{B}

\TabHeading{Atmung}
&
\begin{minipage}[t]{\linewidth}
\begin{itemize}
  -   Schätzen: **Atemfrequenz**, \EE{-tiefe}
  -   ***SpO₂**
  -   Atemgeräusche
  -   Inspektion der \EE{Thorax-Atembewegungen}, ggfs. 
  Atemprobe
  -   *Zeichen der
  Atemnot*: Hautfarbe, Bauch-/Thorax\-bewegungen,
  Anstrengung beim Atmen, Atemhilfsmuskulatur, \ldots
  -   \Z{(Auskultation: Seitenvergleich, Lungenbasen)}
\end{itemize}
\end{minipage}
&
\noindent\color{Indigo}\begin{minipage}[t]{\linewidth}
\begin{itemize}
  -   Reanimation
  -   Assistierte Beatmung (bei AF \textless\,8 oder
  \textgreater\,30\PerMin*)
  -   O₂-Gabe
  -   Situationsgerechte Lagerung
  -   \Z{Erweiterte Maßnahmen\\(Entlastung
  Spannungspneumothorax, \ldots )}
\end{itemize}
\end{minipage}
\\\addlinespace\midrule\addlinespace
\noindent\T{C}

\TabHeading{Kreislauf}
&
\begin{minipage}[t]{\linewidth}
\begin{itemize}
  -   **Hautfarbe**, \EE{-temperatur}
  -   Schweiß
  -   **Rekap**-Zeit
  -   **Radialispuls**: Stärke, Rhythmus, Abschätzen der
  Frequenz
\end{itemize}
\end{minipage}
&
\noindent\color{Indigo}\itshape\begin{minipage}[t]{\linewidth}
\begin{itemize}
  -   Blutstillung
  -   Situationsgerechte Lagerung
\end{itemize}
\end{minipage}
\\\addlinespace\cmidrule{2-3}\addlinespace
\noindent\TabHeading{*STU}
&
\begin{minipage}[t]{\linewidth}
\begin{itemize}
  -   Inspektion
  und Abtasten von
  \begin{inparaenum}
    -   **Kopf** inkl. Ohren,
    -   **Hals**,
    -   **Thorax**,
    -   **Bauch**,
    -   **Becken** und
    -   **Oberschenkel**
  \end{inparaenum}
\end{itemize}
\end{minipage}
&
\noindent\color{Indigo}\begin{minipage}[t]{\linewidth}
\begin{itemize}
  -   Blutstillung
  -   Situationsgerechte Lagerung
\end{itemize}
\end{minipage}
\\\addlinespace\midrule\addlinespace
\noindent\T{D}

\TabHeading{Schnelle\newline Untersuchung}
&
\begin{minipage}[t]{\linewidth}
\begin{itemize}
  -   Vitalwerte (**HF**, **RR**)
  -   Neurologisch:
  \begin{inparaitem}
    -   **Orientierung** oder ***GCS**
    -   **Pupillen**
    -   **Kraftprobe OE**, *UE
    -   *Kann der
    Patient **Hände** und **$2** spüren und aktiv
    bewegen?
  \end{inparaitem}
  -   ***Blutzuckermessung**
\end{itemize}
\end{minipage}
&
\noindent\color{Indigo}\begin{minipage}[t]{\linewidth}
\begin{itemize}
  -   Transportentscheidung (eilig, Rendez-vous, \ldots)
\end{itemize}
\end{minipage}
\\\addlinespace\midrule\addlinespace
\noindent\T{E}

\TabHeading{Erweiterte\newline Untersuchung}
&
\begin{minipage}[t]{\linewidth}
\begin{itemize}
  -   (Fremd-) **Anamnese** (SAMPLE, OPQRST)
  -   Einschätzen der **Umgebung**
  -   Weitere relevante **Untersuchungen**
  -   Arbeitsdiagnose erstellen und Differentialdiagnosen
  ausschließen
\end{itemize}
\end{minipage}
&
\noindent\color{Indigo}\begin{minipage}[t]{\linewidth}
\begin{itemize}
  -   Transportentscheidung 
  -   Spezielle Maßnahmen gemäß Verdachtsdiagnose
\end{itemize}
\end{minipage}
% \\\addlinespace\bottomrule\addlinespace
% \TabHeading{Beurteilung}
% &
% \begin{minipage}[t]{\linewidth}
% \begin{itemize}
%   -   **Einschätzen der vitalen Bedrohung**
%   ständig während des gesamten Blocks
%   -   Vorläufige Verdachtsdiagnose formulieren
% \end{itemize}
% \end{minipage}
% &
% \noindent\color{Indigo}\begin{minipage}[t]{\linewidth}
% \begin{itemize}
%   -   {\TxMassVit}
%   -   Reihenfolge der weiteren Untersuchungen und
%   Maßnahmen festlegen
% \end{itemize}
% \end{minipage}
\\\addlinespace\bottomrule
\end{tabularx}

\vspace{-0.5\baselineskip}\begin{center}
  \EE{Der Einschätzungsblock muss in regelmäßigen Abständen
in angemessenem Umfang wiederholt
werden (Verlaufskontrolle)!}

{\tiny Mit einem * versehene Punkte werden nur durchgeführt,
wenn sie *situationsentsprechend* sind. \enquote{Typische
Sofortmaßnahmen} sind als *Beispiele* zu verstehen.
Bei **absolut zeitkritischen Patienten**, bei denen bereits die 
ABC-Einschätzung ergibt dass ein Transport nicht aufschiebbar ist, 
kann u.\,U. schon \Ca{nach C eine vorgezogene Strategie- 
und Transportentscheidung} getroffen werden. 
D und E sollen dann, sofern möglich, während des 
Transportes erfolgen. }
\end{center}
}
\end{minipage}
\end{minipage}\vspace{-\baselineskip}
}
{}

\vfill
\restoregeometry\pagestyle{fancy}

\Layout{57}{}
{}
{}
{%
\begin{minipage}{\linewidth}
\FontTableBase\footnotesize
{\sffamily\FontTableBase\footnotesize\CorrectionItemizeBox\renewcommand{\itemhook}{%
%   \setlength{\topsep}{0pt}%
  \setlength{\itemsep}{0.005\baselineskip}
%   \setlength{\parskip}{\baselineskip}
  \setlength{\leftmargin}{1.4em}
  \setlength{\labelwidth}{0.9em}
}
\begin{tabulary}{\linewidth}[H]{lCL}
\toprule\addlinespace
{\multirow{6}{2cm}{\TabHeading{Anamnese}}}
 &
{\sffamily\textbf S}
 &
Symptome {\&} Schmerzen: OPQRST

{\tiny(Lokalisation, Stärke,  Qualität, Zeit/Verlauf, 
 Ausstrahlung, Linderung/Verstärkung)}
\\\addlinespace
 &
{\sffamily\textbf A}
 &
Allergien/Allergische Reaktion (auch
Allergien auf Medikamente)
\\\addlinespace
 &
 {\sffamily\textbf M}
 &
Medikamente {\tiny(Nimmt Patient Medikamente?
Wogegen sind diese?)}
\\\addlinespace
 &
{\sffamily\textbf P}
 &
Patientengeschichte: Krankheiten {\tiny(chronische, frühere
K.)}, Operationen, \ldots 
\\\addlinespace
 &
 {\sffamily\textbf L}
 &
Letzte ... (zB letzte Mahlzeit, letzte
Regelblutung, letzter Spitalsaufenthalt)
\\\addlinespace
 &
 {\sffamily\textbf E}
 &
Ereignisse vor Notfalleintritt {\tiny(zB
Anstrengung vor Brustschmerz, Sturz, ...)}
\\\addlinespace\midrule\addlinespace

{\multirow{8}{2cm}{\TabHeading{Diagnostik}}}
 &
{\sffamily AF}
 &
Atemfrequenz auszählen
\\\addlinespace
 &
({\sffamily SpO₂})
 &
Sauerstoffsättigung im Blut, Monitoring
\\\addlinespace
 &
{\sffamily  BZ}
 &
Blutzucker {\tiny(Wann war die letzte
Mahlzeit/Insulinspritze? Diabetes bekannt?)}
\\\addlinespace
 &
{\sffamily Temp.}
 &
Körpertemperatur/Fieber
\\\addlinespace
 &
 {\sffamily\small Neuro\-status}
 &
(Bewusstseinsgrad), Orientierung,
Pupillen, Halbseitenzeichen, Herdblick, Meningismus,
retrograde Amnesie
\\\addlinespace
 &
{\sffamily\small Trauma\-check}
 &
Abnorme Beweglichkeit, DMS an allen
Extremitäten (Nagelbettprobe), Wunden, paradoxe Atmung,
Abwehrspannung des Bauches, ...
\\\addlinespace
 &
{\sffamily\small Körp. Unt.}
&
Sonstige körperliche Untersuchungen: Abdomen abtasten,
Körperstellen inspizieren, \ldots
\\\addlinespace
 &
{\sffamily\small \ldots}
&
EKG, Kapnometrie, \ldots
\\\addlinespace\midrule\addlinespace
\TabHeading{Maßnahmen}
&
&
\begin{minipage}{\linewidth}
\begin{itemize}
  -   (Lagerung)
  -   Verband, Blutstillung
  -   Schienung
  -   psychischer Beistand
  -   \EE{Wärmeerhaltung / Kühlung}
  -   Entscheidung Notarzt-Nachforderung wegen
  Schmerztherapie oder Aufklärung/Belassung
  -   {\TxMassOGabe}
  -   Regelmäßige **Wiederholung des
  Einschätzungsblockes** {\tiny(in angemessenem Umfang)}
  -   **$2**
  (\prettyref{m:standardmassnahmen-immer})
  -   \ldots
\end{itemize}
\end{minipage}
\\\addlinespace\midrule\addlinespace
\TabHeading{Assistenz}
&
&
\begin{minipage}{\linewidth}
\begin{itemize}
  -   EKG anlegen
  -   Infusion/periphervenösen Zugang vorbereiten
  -   Intubation vorbereiten
  -   Medikamente vorbereiten
  -   \ldots
\end{itemize}
\end{minipage}
\\\addlinespace\midrule\addlinespace
\TabHeading{Transport}
&
&
\begin{minipage}{\linewidth}
\begin{itemize}
  -   Welches Spital/Bett buche ich ab? {\tiny(Arbeitsunfall?
  Welche Abteilung fahre ich an? Ist der Patient bereits in
  einem Spital in Behandlung? Spitalsbehandlung vor kurzer
  Zeit (\enquote{Reklamation})?)}
  
  Brauche ich eine Spezialabteilung? {\tiny(Schockraum, Stroke
  Unit, PTCA, Verbrennungsstation)}
  -   Übergabe {\tiny(Anamnese, Befunde, Maßnahmen, Hinweise)}
\end{itemize}
\end{minipage}

\\\addlinespace\bottomrule
\end{tabulary}}

{\scriptsize**Fett** gedruckte Punkte haben eine besonders hohe Priorität, (eingeklammerte) Punkte sollten schon als Teil des Einschätzungsblockes abgehandelt worden sein.}
\end{minipage}
}
{Kurzübersicht: \EEE{Weiterführende Untersuchungen und Anamnese}.
Die Reihenfolge ist je nach Patient
unterschiedlich!\label{tab:checkliste-patientenmanagement}\label{tab:pam-weiterfuehrende}.}
{}
